I encountered clearance issues with the AR0147 BGA fanout. Unlike most BGAs which have a ball radius of 0.5-1.0mm with large pitches, the sensor uses a small package profile. This package made it difficult to properly route the component.
\href{https://resources.altium.com/p/which-bga-pad-and-fanout-strategy-right-your-pcb}{fanout methods} like Via-in-Pad were not possible as such a small scale, simply because the PCB capabilities didn't permit small enough vias.
For most PCB manufacturers, the minimum via diameter size is 0.2mm, with a 0.15mm annular ring (0.35mm total diameter). Since the via was larger than the BGA pad, via-in-pad was not a feasible option.
I then opted for the dogbone fanout method which admit to the manufacturer's trace width and clearance of 3.5mil (see \autoref{fig:fanout}). Unlike the sensor, the STM32N6 has a more coarse BGA package which allowed us to fanout without issues. In this case, I opted for the dog bone fanout since it allowed us to place our ground vias closer to the power vias. It also provided more clearance routing diagonally, allowing us to use thicker traces for a better current carrying capacity.
